
\begin{center}
	\Huge
	Bestemmelse af afledede funktioner
\end{center}

\subsection*{Gruppe 1}

\begin{enumerate}[label=\roman*)]
	\item Vis, at differentialkvotienten for $f(x) = k$ er $f'(x) = 0$ ved at bruge definitionen for
	differentialkvotienten.
	\item Vis, at differentialkvotienten for $g(x) = x$ er $f'(x) = 1$ ved at bruge definitionen for 
	differentialkvotienten
\end{enumerate}

\subsection*{Gruppe 2}
\begin{enumerate}[label=\roman*)]
	\item Vis, at differentialkvotienten for $f$ givet ved
	\begin{align*}
		f(x) = ax + b
	\end{align*}
	er givet ved
	\begin{align*}
		f'(x) = a
	\end{align*}
	ved at bruge definitionen for differentialkvotienten. 
\end{enumerate}

\subsection*{Gruppe 3}
\begin{enumerate}[label=\roman*)]
	\item Vis, at differentialkvotienten for $f$ givet ved
	\begin{align*}
		f(x) = x^3
	\end{align*}
	er givet ved
	\begin{align*}
		f'(x) = 3x^2
	\end{align*}
	ved at bruge definitionen for differentialkvotienten. (Vink: udnyt, at $(x_0+h)^3 = (x_0+h)(x_0+h)^2$ )
\end{enumerate}

\subsection*{Gruppe 4}
\begin{enumerate}[label=\roman*)]
	\item Vis, at differentialkvotienten for $f$ givet ved
	\begin{align*}
		f(x) = ax^2+bx+c
	\end{align*}
	er givet ved
	\begin{align*}
		f'(x) = 2ax+b
	\end{align*}
	ved at bruge definitionen for differentialkvotienten.
\end{enumerate}

\subsection*{Gruppe 5}
\begin{enumerate}[label=\roman*)]
	\item Vis, at differentialkvotienten for $f$ givet ved
	\begin{align*}
		f(x) = \frac{1}{x}
	\end{align*}
	er givet ved
	\begin{align*}
		f'(x) = -\frac{1}{x^2}
	\end{align*}
	ved at bruge definitionen for differentialkvotienten. (Vink: husk reglerne for at trække brøker 
	fra hinanden)
\end{enumerate}

\subsection*{Gruppe 6 (Svær, sorry)}
\begin{enumerate}[label=\roman*)]
	\item Vis, at differentialkvotienten for $f$ givet ved
	\begin{align*}
		f(x) = \sqrt{x}
	\end{align*}
	er givet ved
	\begin{align*}
		f'(x) = \frac{1}{2\sqrt{x}}
	\end{align*}
	ved at bruge definitionen for differentialkvotienten. (Vink: Udnyt, at $(a+b)(a-b) = a^2 - b^2$ hvor du 
	vælger $\sqrt{x_0 + h} -\sqrt{x_0}$ som $(a-b)$ og isolerer dette i udtrykket.)
\end{enumerate}
